\section{Introduction}
\label{sec:intro}

Suppose a server process embeds a SQL engine and lets its users format reports
with \sql{lpad}, which pads a string to a requested width. A user types a width
of 9{,}223{,}372{,}036{,}854{,}775{,}807. The server should answer with an error
message. If the engine instead ends the process, every other user's session ends
with it. A database engine is trusted to be correct without being checked, so a
defect of this kind is caught by the engine's own tests or not at all.

Database engines are among the most heavily tested programs in practice.
Differential testing runs the same query on two engines and reports any
disagreement. Metamorphic oracles compare a query with an equivalent rewrite of
itself~\cite{rigger2020tlp,rigger2020norec}. Grammar-based generators produce
valid SQL at scale~\cite{zhong2020squirrel,liang2026sqlbull}. All of these tools
stop at two boundaries: the generator can only emit the functions and values it
was written to emit, and a differential oracle can only judge a feature that its
reference engine also implements.

We focus on the defects that sit past those two boundaries. Such defects are
silent in a specific way: a fuzzer can run for a year without reaching them, and
the time it has already run says nothing about how close it is.

\begin{figure}[t]
\begin{lstlisting}[escapechar=@]
-- (a) Regression test added by the fix for DuckDB issue #27846
select LPAD('x', 10, '0'), LPAD('x', 4, '@\textcolor{black}{\ding{72}}@'),
       LPAD('x', 3, '@\textcolor{black}{\ding{72}}@ab'), RPAD('x', 3, '@\textcolor{black}{\ding{72}}@ab')
-- (b) The same shape with every literal replaced by a typed hole
select LPAD({s4}, {n0}, {s5}), LPAD({s6}, {n1}, {s7}),
       LPAD({s8}, {n2}, {s9}), RPAD({s10}, {n3}, {s11})
-- (c) Filled from a boundary pool at seed 0, then run on Turso b588fbdc6
select LPAD('', -9223372036854775808, 'abc'), LPAD('_', -128, '%'),
       LPAD('0', 9223372036854775807, 'abc'), RPAD('%', -128, 'ABC')
-- (d) Reduced to one statement
select LPAD('a', 9223372036854775807, 'b');
@\textit{memory allocation of 9223372036854775806 bytes failed}@      @$\leftarrow$ process aborts, exit status 134@
\end{lstlisting}
\caption{The motivating example. A regression test from another engine~(a),
templated~(b) and refilled with boundary values~(c), aborts Turso from one
statement~(d). The marker \ding{72} stands for a four-byte emoji in the original
test. Turso's own fuzzers cannot produce~(d): their 47-function vocabulary has
no \sql{lpad}, and their integers stop at 1{,}000{,}000.}
\label{fig:motivating}
\end{figure}

Figure~\ref{fig:motivating} shows a real case. Turso~\cite{turso} is a
SQLite-compatible engine written in Rust. It runs its own differential fuzzer in
continuous integration: 41{,}205 lines of Rust that generate SQL and compare
Turso's answers with SQLite's. Yet the statement in Figure~\ref{fig:motivating}d
ended the process with exit status 134 (SIGABRT) in 2.28 seconds, where it should
have returned the engine's own error \emph{string or blob too big}. Turso's
fuzzer could never have produced it. Its vocabulary of 47 scalar functions does
not contain \sql{lpad}, and it draws integer literals only from
$[-1{,}000{,}000,\ 1{,}000{,}000]$~\cite{tursofuzzer}. Even if it had produced
the statement, SQLite has no \sql{lpad} to compare the answer against. This is
not a single case: we found four such defects, and none of Turso's four
libFuzzer targets or its differential fuzzer calls any of the three functions
involved.

The closest prior work either chooses where to look from the grammar or reuses
known defects through static analysis. SQLBull starts generation from grammar
rules that a tester marks as feature-rich~\cite{liang2026sqlbull}.
Recurring-vulnerability detectors turn a known patch into a rule and search for
code that breaks it~\cite{cao2025vulpa,cao2025antman}. However, the first needs
someone to pick the starting rules, and the second needs the defect to recur in
shared or copied code. Turso and DuckDB share no code.

We present \tool, which finds defects in this region in four steps:
\begin{enumerate}
  \item It reads the incumbent fuzzer's source and lists the functions and
        values that fuzzer cannot generate (Section~\ref{sec:step-blind}).
  \item It mines closed, reproduced bug reports from a different engine and keeps
        each report's query as a template whose literals are typed holes
        (Section~\ref{sec:step-mine}).
  \item It fills those holes, and the output of a typed SQL generator, with
        boundary values, and runs every statement against the engine under test
        (Section~\ref{sec:step-run}).
  \item It flags any statement that ends the process, then reduces, re-checks and
        reports it (Section~\ref{sec:step-triage}).
\end{enumerate}

We ran \tool against Turso. It found four previously unknown defects in Turso's
string functions: \sql{lpad} and \sql{rpad} allocate their result before checking
its length, \sql{repeat} overflows its size computation, and the length check
counts characters against a limit in bytes. Our two patches, which fix all four,
were merged upstream~\cite{turso11225,turso11259}. At an equal budget of
20{,}000 cases, the typed generator reached three distinct defects and the mined
reports reached two. The lesson we draw is a practical one: before spending
compute on a fuzzer, spend an hour reading the fuzzer that is already there.

\paragraph{Contributions.}
\begin{itemize}
  \item We define the \emph{blind spot} of an incumbent fuzzer and show that it
        can be measured from that fuzzer's source (Section~\ref{sec:problem}).
  \item We present \tool, which fills the blind spot with the shapes of another
        engine's bug reports and with typed generation, and judges results with an
        oracle that needs no reference engine (Section~\ref{sec:approach}).
  \item We report four previously unknown defects in Turso, all fixed upstream by
        our patches, each regenerated from a recorded seed
        (Section~\ref{sec:eval}).
  \item We measure what each input source contributes at an equal budget, and
        release the tool, the corpus and all data
        (Sections~\ref{sec:rq2} and~\ref{sec:data}).
\end{itemize}
